\section{Proofs for the complete paid-support calculation}
\label{supp:memory123}

\subsection{Derivative transport, integration and the optimal upper-model proposal}

\begin{proof}[Proof of Theorem~\ref{thm:memorycontrast123}]
For $s<t$, differentiable strong convexity implies
$q'(t)-q'(s)\ge\mu(t-s)$: apply its first-order inequality in both
directions and add. Derivative Lipschitzness gives
$q'(t)-q'(s)\le L(t-s)$. At an anchor $a_i$, transporting the
lower endpoint to the right therefore uses $\mu$, whereas transporting
it to the left uses $L$. The upper endpoint uses $L$ to the right
and $\mu$ to the left. This gives
$\lambda_i(t)\le q'(t)\le\upsilon_i(t)$ and hence the first
inequality in \eqref{eq:memoryvalid123}. Integrating it over the
oriented interval from $a$ to $b$ proves its second inequality,
including the reversal for a negative displacement.

For $b\ge a$, pointwise
$\upsilon\le\upsilon_i$ and $\lambda\ge\lambda_i$ give
$C^+\le C_i^+$ and $C^-\ge C_i^-$. For $b<a$, the integrals
acquire a minus sign and the two envelope roles exchange; the same
two conclusions follow. Taking the extrema over all $i$ proves
\eqref{eq:memorydominance123}. Adding a sound anchor increases
$\lambda$ and decreases $\upsilon$ pointwise, which proves
monotonicity of the contrast bounds. It is not necessary to assume
that the new anchor was an accepted policy.

Every function $\lambda_i$ and $\upsilon_i$ is continuous. For
$s<t$ its increment lies in $[\mu(t-s),L(t-s)]$, also when its
kink lies between $s$ and $t$. Maxima and minima preserve these
increment inequalities. For example, an index maximizing
$\lambda_i(s)$ gives the lower increment inequality for
$\lambda$, while an index maximizing $\lambda_i(t)$ gives its
upper inequality. For a minimum choose the minimizing index at
$t$ for the lower inequality and at $s$ for the upper inequality.
Thus both global envelopes are continuous, strictly increasing,
and $L$-Lipschitz.

For $b<a$, the derivative of $C^+(a,b)$ in $b$ is
$\lambda(b)$; for $b>a$ it is $\upsilon(b)$. The one-sided
derivatives at $a$ are $\lambda(a)$ and $\upsilon(a)$, in that
order, and are ordered because they enclose $q'(a)$. On each
side subtracting $\mu b$ from these derivatives produces a
nondecreasing function, and their jump at $a$ is nonnegative.
Integration proves that $C^+(a,b)-\mu b^2/2$ is convex.
Therefore $C^+$ is $\mu$-strongly convex and has a unique minimizer
on the compact interval. The subgradient condition at $a$, the
strict increase of both envelopes, and the constrained endpoint
conditions give exactly the stated three cases. If an endpoint
derivative equals zero, that endpoint is the unique minimizer.
The proof also covers $a=0$, $a=A$, repeated anchors and $L=\mu$.
\end{proof}

The theorem does not infer a derivative bound on the missing side of
a boundary support. It starts only after every anchor interval has
been validly obtained, using the boundary convention of
Theorem~\ref{thm:supportgradient122} or other proved original
information. If a missing-side bound is infinite, it may be retained
in the relevant envelope algebra, but the finite rational algorithm
below presumes finite supplied intervals. The original investment
gradient cap supplies finite intervals where applicable.

\subsection{Exact finite arithmetic and robust lower supports}

Suppose the anchors, all interval endpoints, $A$, $\mu$ and $L$
are rational. Split $[0,A]$ first at every anchor. On each resulting
interval, each local derivative bound is affine. Include every
intersection of every pair of these affine functions which lies in
that interval. There are finitely many, and all are rational; an
identical pair contributes no new point. Between consecutive
partition points, a rational midpoint determines an active maximum
and minimum. Any fixed index order resolves an exact tie without
changing its value. Integrating the active affine expression gives
the exact envelope integral over that piece. A zero in an active
piece is rational because its slope is at least $\mu>0$.

Equivalently, retain the at most two affine pieces from each anchor,
each with its applicable interval, and generate only pairwise
intersections whose piece domains overlap. This gives at most a
quadratic number of candidate breakpoints in $m$; sorting and testing
all candidates is one finite algorithm. No particular optimized
operation count is claimed for an implementation that instead
loops over all pieces on all intervals. It must record its actual
arithmetic and time costs.

At $L=\mu$, every derivative piece has the same slope. The envelopes
are the largest and smallest corresponding intercepts plus $\mu t$.
There is no division by $L-\mu$. At $A=0$ the only feasible action is
retention and every contrast is exactly zero; the nontrivial
statements above are used only for $A>0$.

For the robust support calculation, fix a piece on one side of every
anchor. Subtract $\mu b^2/2$ from \eqref{eq:rawsupport123}. The result
is affine on that piece:
\[
 s_i(b)-\frac\mu2 b^2
   =(p_i^\diamond-\mu a_i)b
     +W_i-p_i^\diamond a_i+\frac\mu2a_i^2,
\]
where $p_i^\diamond=p_i^+$ on $b<a_i$ and
$p_i^\diamond=p_i^-$ on $b>a_i$. Add the pairwise intersections
of these affine functions within their applicable intervals.
On each resulting piece, $\max_i s_i(b)$ is a single quadratic
with curvature $\mu$. Evaluate it at both endpoints and at its
stationary point if that point lies inside the piece. The least of
these finitely many rational values is the exact global minimum.
The endpoint evaluations include every downward derivative jump.
This procedure proves existence and computes the required minimum
without incorrectly calling the whole robust envelope convex.

With exact slopes, $p_i^-=p_i^+$, the robust supports really are
convex quadratics with common curvature and their maximum is
$\mu$-strongly convex. That useful special case does not justify
the same assertion for interval slopes.

\begin{proof}[Proof of Theorem~\ref{thm:memorypolicy123}]
The support hypothesis says that there is an actual
$p_i\in[p_i^-,p_i^+]$ with
\[
 q(b)\ge W_i+p_i(b-a_i)+\mu(b-a_i)^2/2.
\]
Minimizing the linear term over its interval gives
$q(b)\ge s_i(b)$ for every feasible $b$. Equation
\eqref{eq:memoryvalid123}, together with the value bounds at $a_i$,
also gives
$W_i+C^-(a_i,b)\le q(b)\le U_i+C^+(a_i,b)$.
Taking the indicated maximum and minimum proves
$\underline q\le q\le\overline q$. Minimizing the global
support inequality gives
$\min_b\max_i s_i(b)\le\min_b q(b)$, and taking a maximum
with another valid optimum lower bound preserves that order.

For distinct actions, $\eta(b)-\eta(a)\in[-\delta,\delta]$.
Adding this interval to the original $q$ contrast proves the first
line of \eqref{eq:memorypolicy123}. Since $0\le\eta(b)\le\delta$,
the absolute bounds on $q(b)$ imply the stated bounds on $J(b)$.
The original implemented future is feasible and has nonnegative
error, so $J(b)\ge q(b)\ge\min q$. Subtracting the lower optimum
bound proves the regret upper bound. Subtracting the absolute
policy bounds at $a,b$ gives the additional signed interval; both
contain the same true contrast, so their intersection is valid.

For $b=a$, the actually implemented action and future policy coincide
and subtraction is identically zero. An allowance $\delta$ is
not charged to this identity. For an actual change, the three
acceptance inequalities respectively certify strict original-policy
improvement, original regret qualification, and original objective
width qualification. Rounding an ideal proposal changes the
argument of every bound; the proof applies to its actual rounded
feasible value because the bounds hold for every such value.
For a finite accepted sequence, summing the original-policy contrast
inequalities telescopes. This proves all statements.
\end{proof}

If a prospective implementation adds anchors after observing earlier
certificate outcomes, these deterministic inequalities remain valid
for every sound realized collection. They do not establish a
probabilistic claim for adaptively acquired labels; a stochastic
oracle needs simultaneous coverage for every retained adaptive
query. Nor does a retrospective reanalysis become a prospectively
frozen experiment by applying deterministic bounds to it.

\subsection{Sharpness at the explicitly reduced derivative interface}

Fix $A>0$, $0<\mu\le L$, and finitely many pairs
$(a_i,[\ell_i,u_i])$. Let $\mathcal Q_D$ be the set of all
continuously differentiable, $\mu$-strongly convex functions on
$[0,A]$ with $L$-Lipschitz derivative satisfying all these derivative
intervals. No value, support, gradient-cap or investment-formula
information beyond these intervals is imposed. Assume
$\mathcal Q_D\ne\varnothing$. Fix a feasible incumbent $a$.

\begin{proposition}[Sharp deterministic and randomized gain guarantees]
\label{prop:memoryminimax123}
For each feasible $b$,
\[
 \sup_{q\in\mathcal Q_D}\{q(b)-q(a)\}=C^+(a,b),\qquad
 \inf_{q\in\mathcal Q_D}\{q(b)-q(a)\}=C^-(a,b).
\]
Both extrema are attained. Put
$\Gamma=-\min_{b\in[0,A]} C^+(a,b)\ge0$.
An admissible method receives only this fixed information and
chooses any Borel probability law $\nu$ for $B\in[0,A]$ without
new queries. The adverse function knows $\nu$ before $B$ is drawn.
Then
\[
 \inf_\nu\sup_{q\in\mathcal Q_D}
      \mathbb E_\nu[q(B)-q(a)]=-\Gamma.
\]
If $J=q+\eta$ and $\eta$ ranges over all Borel functions from
$[0,A]$ to $[0,\delta]$, then
\[
 \inf_\nu\sup_{q,\eta}
      \mathbb E_\nu[J(B)-J(a)]
       =\min\{0,-\Gamma+\delta\}.
\]
Retention or the deterministic minimizer of $C^+(a,\cdot)$ attains
the latter value. This is the exact guarantee
\eqref{eq:memoryminimax123} at the stated interface.
\end{proposition}

\begin{proof}
Choose any $q_0\in\mathcal Q_D$. The derivative-envelope proof
gives $\lambda\le q_0'\le\upsilon$ pointwise. At every $a_i$,
$\lambda(a_i)\ge\ell_i$ and
$\upsilon(a_i)\le u_i$ by the terms carrying that very anchor.
Together these facts imply
\[
 \ell_i\le\lambda(a_i)\le u_i,
 \qquad \ell_i\le\upsilon(a_i)\le u_i.
\]
The preceding increment proof shows that each envelope is
continuous and has all secants in $[\mu,L]$. Its indefinite
integral is therefore continuously differentiable, has
$L$-Lipschitz derivative, and is $\mu$-strongly convex. It also
satisfies every anchor interval. Thus, for any constant $c$,
\[
 q_\lambda(t)=c+\int_a^t\lambda(s)\,ds,
 \qquad
 q_\upsilon(t)=c+\int_a^t\upsilon(s)\,ds
\]
belong to $\mathcal Q_D$. If $b\ge a$, the second function
attains $C^+$ and the first $C^-$; if $b<a$, their roles exchange.
The general bounds were already proved, so this establishes
both sharp equalities without asserting independent choices of
a worst derivative at each point.

For any law $\nu$, convexity of each admissible $q$ and Jensen's
inequality give
\[
 \mathbb E_\nu[q(B)-q(a)]
   \ge q(\mathbb E_\nu B)-q(a).
\]
All expectations are finite because $q$ is continuous on a compact
interval. Taking the supremum and using the just-proved sharp
fixed-action bound gives a value at least
$C^+(a,\mathbb E_\nu B)\ge-\Gamma$. The deterministic minimizer
attains $-\Gamma$, proving the first minimax equality.

For the implemented error put $r=\Pr_\nu(B\ne a)$. If $r=0$,
the contrast is exactly zero. If $r>0$, set
$b_*=\mathbb E_\nu[B\mid B\ne a]\in[0,A]$. Conditional Jensen
gives, for every admissible $q$,
\[
 \mathbb E_\nu[q(B)-q(a)]
 \ge r\{q(b_*)-q(a)\}.
\]
The single Borel function which is zero at $a$ and $\delta$
elsewhere gives expected error difference $\delta r$.
It may be chosen together with a function attaining the
fixed-action supremum at $b_*$. Consequently
\[
 \sup_{q,\eta}\mathbb E_\nu[J(B)-J(a)]
 \ge r\{C^+(a,b_*)+\delta\}
 \ge r(-\Gamma+\delta)
 \ge\min\{0,-\Gamma+\delta\}.
\]
This also covers $b_*=a$ when the conditional law straddles the
incumbent. Retention attains zero. If $\Gamma>\delta$, its
positive value implies that the minimizing action differs from
$a$, and the sharp fixed-action contrast plus an error difference
$\delta$ attains $-\Gamma+\delta$. At equality both choices are
optimal when distinct, and the zero case is covered by retention.
No adversary observes the realized random action. This proves
the second equality.
\end{proof}

This proposition is deliberately stated for $\mathcal Q_D$.
Integrating the extremal derivative may fail additional absolute
value constraints or the algebraic restrictions of the investment
model. Soundness and dominance of the envelope calculation persist
after those restrictions, but its sharpness and minimax claims
are not transferred to that narrower class. No numerical adoption
or representation-specific conclusion follows from the proposition.

\subsection{An exact rational strict-information witness}
\label{supp:memorywitness123}

Let $A=3/2$, $\mu=1$, $L=4$, incumbent $a=1/4$, and derivative
anchors
\[
 a_1=1/4,\quad \ell_1=u_1=-2,
 \qquad a_2=5/4,\quad\ell_2=u_2=-1/2.
\]
Write $x=t-1/4$. Between the two anchors the upper envelope is
\[
 \upsilon(t)=\min\{-2+4x,\ x-3/2\}.
\]
The first term is active for $0\le x\le1/6$ and the second
for $1/6\le x\le1$. To the right of the second anchor,
$\upsilon(t)=4x-9/2$. Its zero is $x=9/8$, hence the joint
robust proposal is $b=11/8$. Its exact upper contrast is
\begin{align*}
 C^+(a,11/8)
 &=\int_0^{1/6}(-2+4x)\,dx
   +\int_{1/6}^{1}(x-3/2)\,dx
   +\int_1^{9/8}(4x-9/2)\,dx\\
 &=-\frac5{18}-\frac{55}{72}-\frac1{32}
  =-\frac{103}{96}.
\end{align*}
Using only the first anchor gives minimizer $x=1/2$ and guaranteed
gain $1/2$. Using only the second gives minimizer $x=9/8$ and
upper contrast
\[
 \int_0^1(x-3/2)\,dx+
 \int_1^{9/8}(4x-9/2)\,dx=-\frac{33}{32}.
\]
The joint guaranteed gain therefore exceeds the better individual
guarantee by $103/96-33/32=1/24$. The admissible function
$q(t)=\int_{1/4}^t\upsilon(s)\,ds$ attains the joint upper
contrast. Its derivative is continuous with slopes either one
or four, and it satisfies both exact anchor derivatives.
Thus the example verifies a realizable smooth-convex information
improvement, not an intersection of mutually incompatible derivatives.

At the two anchors this witness has levels $q(a_1)=0$ and
$q(a_2)=-25/24$. Exact touching $\mu$-quadratic supports at these
interior anchors are valid for this function. If those exact
levels are additionally disclosed, an algorithm can exploit the
known cross-anchor value difference as well; the reduced-interface
lower bound is not a lower bound for that stronger information
set. The example is therefore not advertised as proving a strict
advantage over every possible full-payload algorithm.

The example is algebraic validation. It is not an experimental
state from the investment application, a new population sample,
or evidence of a uniquely neural advantage. An independently
implemented rational calculation can reproduce every number
without importing an economic or envelope production module.

\subsection{Optimizer brackets and paid action-directed precision}

\begin{proof}[Proof of \eqref{eq:bracket123} and Proposition~\ref{prop:midpoint123}]
The continuously differentiable strongly convex function $q$ has a
unique minimum on the compact interval. Its constrained first-order
conditions give $a_q^*=z(q')$, with the endpoint conventions stated
in the main text. If two continuous strictly increasing functions
satisfy $g\le h$, then $z(h)\le z(g)$. To verify this directly,
consider an interior zero of either function and use their sign
on either side; if either constrained zero is an endpoint the
same sign argument gives the weak inequality. Applying this
order property to $\lambda\le q'\le\upsilon$ proves
\eqref{eq:bracket123}.

An additional anchor increases $\lambda$ and decreases
$\upsilon$, so $z(\lambda)$ can only decrease and
$z(\upsilon)$ can only increase. Thus the new bracket is a
subset of the old one. This gives the first bound $w_{\rm new}\le w$.

Write the newly purchased derivative interval as $[d_-,d_+]$,
with $d_+-d_-\le\epsilon$. If $d_-\ge0$, the new lower
envelope is nonnegative at $c$, so its constrained zero is at most
$c$. The new bracket is contained in
$[\underline a,\min\{\overline a,c\}]$ and has width at most
$w/2+r$. If $d_+\le0$, its upper envelope is nonpositive at $c$;
the same argument gives containment in
$[\max\{\underline a,c\},\overline a]$ and the same width bound.
These assertions also hold when the new interval has a zero
endpoint, because $\mu>0$ makes both envelopes strictly increasing.

In the remaining case $d_-<0<d_+$, the new upper anchor cone
on the left of $c$ is $d_++\mu(t-c)$. Its constrained zero
cannot be less than $c-d_+/\mu$, clipped at zero. The new
lower anchor cone on the right is $d_-+\mu(t-c)$, whose
constrained zero cannot exceed $c-d_-/\mu$, clipped at $A$.
Since the fused envelope is at least as informative as this
single new interval, its bracket is contained in
\[
 [c-d_+/\mu,\ c-d_-/\mu]\cap[0,A].
\]
This has width at most $(d_+-d_-)/\mu\le\epsilon/\mu$.
Combining the three cases and the nested-bracket bound proves
\eqref{eq:bracketcontraction123}. It makes no assumption about
which sign the purchased interval will have. Applying the stated
relative precision and rounding conditions recursively proves
$w_k\le\gamma^k w_0$.

For the potential ceiling, the true optimizer belongs to the
bracket and the lower contrast bound gives
$q(a)-q(a_q^*)\le-C^-(a,a_q^*)$. Maximizing over the entire
bracket and taking the maximum with zero proves
\eqref{eq:bracketpotential123}. For every other feasible $b$,
$J(a)-J(b)\le q(a)-\min q+\delta\le K_q(a)+\delta$.
The original incumbent's regret bound is another valid ceiling;
their minimum is therefore valid. The same proof as
Proposition~\ref{prop:pricedstop122} establishes the stated
single-decision cost dominance when a compulsory incremental
query price exceeds that ceiling.
\end{proof}

The midpoint does not have to be an unrounded real number. The
geometric loss caused by its actual feasible representation is
the explicit $r$ term. If the prescribed midpoint equals a paid
anchor and the protocol declines to buy it again, no contraction
is asserted. If the prescribed query returns a broad interval,
the exact new bracket may be unchanged. Both cases have their
original status and cost rather than being counted as successful
action resolution.
